\documentclass[11pt]{article}
\usepackage[margin=1in]{geometry}
\usepackage{amsmath,amssymb,amsthm,mathtools}
\usepackage[hidelinks]{hyperref}
\newtheorem{theorem}{Theorem}
\newtheorem{lemma}[theorem]{Lemma}
\newtheorem{corollary}[theorem]{Corollary}
\newtheorem{proposition}[theorem]{Proposition}
\theoremstyle{remark}
\newtheorem{remark}[theorem]{Remark}
\newcommand{\Had}{\mathcal H}
\newcommand{\hor}{\mathrm h}
\newcommand{\e}{\mathrm e}
% Step 4 numbers (filled in from proof/atlas/results/)
\newcommand{\NRecords}{$12{,}021$}
\newcommand{\Fsix}{$877$}

\title{There are no four mutually unbiased bases in $\mathbb C^6$:\\ a four-step proof}
\author{}
\date{8 October 2026}

\begin{document}
\maketitle

\begin{abstract}
We prove that $\mathbb C^6$ admits no four mutually unbiased bases. The argument has four steps.
\begin{enumerate}
\item An exact algebraic fact about a four-variable seed polynomial $R$: its positive local maxima form a single orbit, with $\max R=3^{13}/2^{12}$.
\item A topological consequence, relative to the classification of order-6 complex Hadamard matrices: the space of non-Tao Hadamard classes is connected.
\item A local exclusion lemma for balls around arbitrary Hadamard matrices.
\item A finite certified atlas of such balls whose boundaries transfer into one another.
\end{enumerate}
Steps 1 and 2 are pure mathematics. Steps 3 and 4 are certified computation in interval arithmetic.
\end{abstract}

\section{Setting}

\paragraph{Hadamard matrices.} Let $\Had$ denote the order-6 complex Hadamard matrices: $|H_{ij}|=1$ and $HH^*=6I$. Write $H\cong H'$ if $H'=PD_1HD_2Q$, with $P$ and $Q$ permutation matrices and $D_1$ and $D_2$ diagonal unitaries.

\paragraph{Quartet edges.} After a unitary change of frame, four mutually unbiased bases (MUBs) contain $I$ and $H/\sqrt6$ for some $H\in\Had$. Call $H$ a \emph{quartet edge} if $\{I,H/\sqrt6\}$ extends to four MUBs. This property is invariant under:
\begin{itemize}
\item equivalence;
\item conjugation $H\mapsto\bar H$;
\item transposition, since applying $H^*/\sqrt6$ exchanges the roles of $I$ and $H/\sqrt6$.
\end{itemize}
We show that no $H\in\Had$ is a quartet edge.

\paragraph{The quotient metric.} A real $6\times6$ matrix is \emph{horizontal} if all its row and column sums vanish. Let $\hor$ be the orthogonal projector onto horizontal matrices. Its kernel consists of the row-plus-column patterns, so $\e^{iY}\cong\e^{i\hor(Y)}$, with exponentials taken entrywise. On $\Had/{\cong}$ we use
\[
 \delta(H,K)=\min\{\|v\|:\ v\text{ horizontal},\ H\cong K\circ \e^{iv}\}.
\]
Let $S_6$ be Tao's matrix and $X=(\Had/{\cong})\setminus\{[S_6]\}$.

\paragraph{Corner seeds.} For $\theta\in\mathbb T^4$, let $E(\theta)=\begin{psmallmatrix}1&1&1\\1&a&b\\1&c&d\end{psmallmatrix}$ with $(a,b,c,d)=\e^{i\theta}$, and $G=6I-EE^*$. Set $s=G_{21}$, $t=G_{31}$ and $r=G_{32}$, and form the invariants $g,\eta,\chi$ of the classification. Then
\[
 R(\theta)=4g^2\eta-\chi^2 .
\]
On $S=\{R>0\}$ an explicit formula produces a Hadamard completion $H_+(\theta)$ of $E(\theta)$. Relative to the classification, statements MC1--MC5 give
\[
 X=\{[H_+(\theta)]\}\cup\{[\overline{H_+(\theta)}]\}\cup\{[F_6]\}\qquad(\theta\in S).
\]

\section{Step 1: the maxima of $R$}

\begin{theorem}\label{thm:max}
The positive local maxima of $R$ on $\mathbb T^4$ are exactly the $18$-point orbit, under the $144$ seed symmetries, of
\[
 \theta_*=\alpha(1,2,-1,-2),\qquad \alpha=\arccos\tfrac14,\qquad R(\theta_*)=3^{13}/2^{12}.
\]
The Hessian at $\theta_*$ is negative definite.
\end{theorem}

\begin{proof}[Outline]
\emph{The deltoid traces.} Normalise the even products $(ad,b,c)$ and the odd products $(bc,d,a)$ to have product $1$, and let $U$ and $V$ be their traces. Each lies in the filled deltoid $\{\,Z:|Z|^4+18|Z|^2-8\operatorname{Re}Z^3-27\le0\,\}$, and the trace map onto it is open.

\emph{The compact form of $R$.} With $A=U+V$, $B=U-V$, $u=|A|^2$ and $v=|B|^2$, one has $g=2v-3u$, $\eta=v-2u$ and
\[
 \chi=\tfrac12(3u-v)(v-u+12)-2\operatorname{Re}(AB^2).
\]
Thus $R$ is invariant under the circle action $(A,B)\mapsto(\e^{-2i\tau}A,\e^{i\tau}B)$.

\emph{No interior maxima.} Every critical point of $\mathcal R(A,B)$ on $\mathbb C^2$ with $\mathcal R>0$ has a direction of positive second derivative. Rotate to $B=\sqrt v>0$ and write $A=x+iy$; positivity gives $d=v-2u>0$. Then:
\begin{itemize}
\item $H_v=4g(6v-11u)-2CC_v$ with $g,\,6v-11u>0$, so $C\neq0$;
\item $yH_x-xH_y=4vyC$, so $y=0$;
\item at $x=-h\le0$ and $v=2h^2+d$, the Jacobian $c_xf_v-c_vf_x$ equals $-8(h^2+2d)W$ with $W\ge6d^2>0$, so $x>0$;
\item $f_v=0$ then gives $c<0$, and $xH_{yy}=-4vc>0$.
\end{itemize}

\emph{Both traces on the boundary.} A smooth boundary trace with an interior partner moves inward along the circle action, by the curvature identity $-\tfrac92(r+2)^2|z^3-1|^4<0$. A cusp with an interior partner is excluded by $\mathcal R=-4FG^2$ and a Bernstein certificate.

\emph{The double boundary.} Here $\mathcal R=-256L^2M^2N$. An explicit eliminant, with cofactors of $11$ and $22$ terms, leaves one stationary branch. On it, the transverse curvature exceeds $4125/32$. Bernstein certificates on $h=1$ leave only $(c,h)=(1/4,1)$, whose lifts form the orbit.

There are $651$ exact claims, checked with integers and rationals only. Step~1 is proved in Lean in full: the $216$ certificate theorems together with the analytic assembly, using only the standard axioms. No resultant of high degree, Gr\"obner basis or Sturm count is needed. Two earlier, independent proofs remain as cross-checks.
\end{proof}

\section{Step 2: connectedness}

\begin{theorem}\label{thm:conn}
Relative to MC1--MC5, the space $X$ is connected.
\end{theorem}
\noindent This upgrades the 2007 observation of Bengtsson et al.\ \cite{BBELTZ}, that the then-known order-6 families form a connected set apart from $S_6$, to a theorem about all order-6 Hadamard matrices.

\begin{proof}
\emph{Sheet parity.} If a seed symmetry $L$ has transposition bit $\tau$ and conjugation bit $\kappa$, then
\[
 [H_\sigma(L\theta)]=\mathcal C^\kappa\mathcal T^\tau[H_{(-1)^{\tau+\kappa}\sigma}(\theta)].
\]
The maximal frame $E(x)=\begin{psmallmatrix}1&1&1\\1&x&x^2\\1&x^{-1}&x^{-2}\end{psmallmatrix}$ has two exact symmetries. Together these give $H_+(L\theta_*)\cong K_*$ for all $144$ maps $L$.

\emph{Every component reaches $[K_*]$.} The class map $\theta\mapsto[H_+(\theta)]$ is continuous on $S$, including across root cuts and sign switches. Every component of $S$ has compact closure, with $R>0$ inside and $R=0$ on its boundary, so it contains a local maximum of $R$. By Theorem~\ref{thm:max} that maximum lies in the orbit of $\theta_*$. So $A=\{[H_+(\theta)]\}$ is connected and contains $[K_*]$, and likewise $\bar A$.

\emph{Joining at $F_6$.} The family $\begin{psmallmatrix}F_3&F_3\\F_3D(t)&-F_3D(t)\end{psmallmatrix}$, with $D(t)=\operatorname{diag}(1,\e^{it},1)$, tends to $F_2\otimes F_3\cong F_6$. Its cross-ratio $\e^{it}$ keeps it away from $F_6$ and $S_6$ for $0<t<\pi/3$. So $[F_6]$ lies in the closure of both $A$ and $\bar A$, and $X=A\cup\bar A\cup\{[F_6]\}$ is connected.
\end{proof}

\begin{corollary}\label{cor:cover}
Let $\mathcal R$ be a finite family of closed balls $B_r=\{\e^{i(P_r+\delta)}:\delta \text{ horizontal},\ \|\delta\|\le R_r\}$. Let $W$ be the set of classes of Hadamards in the union of its closure under $\Gamma=\{\mathrm{id},\text{conjugation},\text{transposition},\text{conjugate-transpose}\}$. Suppose:
\begin{enumerate}
\item $W$ contains a non-Tao Hadamard;
\item every Hadamard with $\|\delta\|=R_r$, for $r\in\mathcal R$, is equivalent after some $g\in\Gamma$ to one with $\|\delta'\|<R_t$, for some $t\in\mathcal R$.
\end{enumerate}
Then every class in $X$ meets $W$.
\end{corollary}

\begin{proof}
$W\cap X$ is closed, as a finite union of compact images. It is open, since every point has an interior representative and $\delta$ is invariant under equivalence. It is nonempty. Since $X$ is connected, $W\cap X=X$.
\end{proof}

\section{Step 3: local exclusion}

Let $\Phi(\delta)\in\mathbb R^{30}$ be the off-diagonal entries of $(\e^{i(P+\delta)})(\e^{i(P+\delta)})^*$, and $J=\Phi'(0)$. Entry $(a,b)$ of $\Phi(\delta)-\Phi(0)-J\delta$ is $\sum_c Z_{ac}\bar Z_{bc}(\e^{iy_c}-1-iy_c)$ with $y_c=\delta_{ac}-\delta_{bc}$. Since $|\e^{iy}-1-iy|\le y^2/2$, $\|\delta_a-\delta_b\|^2\le2\|\delta\|^2$ and $\sum_{a<b}\|\delta_a-\delta_b\|^2\le6\|\delta\|^2$, we get, for every real $\delta$ and every centre,
\[
 \|\Phi(\delta)-\Phi(0)-J\delta\|\le\tfrac L2\|\delta\|^2,\qquad L=2\sqrt3 .
\] A \emph{ball Record} is a rational centre $P$, a rational tangent frame $T$ ($36\times k$), and a constant $\sigma>0$ with $\|J\nu\|\ge\sigma\|\nu\|$ on the horizontal complement $N$ of $T$. The constant $\sigma$ is certified by an Arb Cholesky factorisation.

\begin{lemma}[Normal bound]\label{lem:normal}
Every Hadamard $\e^{i(P+\delta)}$ with $\|\delta\|\le R<\sigma/L$ has $\delta=T\eta+\nu$ with $\|\eta\|\le r_H$ and $\|\nu\|\le n$, where
\[
 n=\frac{\sigma-\sqrt{\sigma^2-2Lb}}{L},\qquad b=\|\Phi(0)\|+\|JT\|r_H+\tfrac L2\lambda_{\max}(T^\top T)\,r_H^2 .
\]
\end{lemma}

\begin{proof}
We have $0=\Phi(\delta)=\Phi(0)+J\delta+E$ with $\|E\|\le\frac L2\|\delta\|^2$, and $\|\delta\|^2=\|T\eta\|^2+\|\nu\|^2$. Hence $\sigma\|\nu\|-\frac L2\|\nu\|^2\le b$. Since $\|\nu\|\le R<\sigma/L$, $\|\nu\|$ lies below the smaller root.
\end{proof}

\begin{lemma}[Stage C, after {[OAI, \S6]}]\label{lem:stagec}
Let $r_{\rm col}$ bound the columns of $\delta$ and $r_F$ its Frobenius norm. Suppose the vector-bin cover, the label clustering, the O/U edge tests (including the joint test, which uses the common $\eta$) and the exhaustive two-clique search succeed with these bounds. Then no Hadamard in the ball is a quartet edge.
\end{lemma}

All inequalities are proved in \texttt{proof/atlas/STAGEC.md}. Comparisons use Arb-tabulated exponentials and float arithmetic with margin $10^{-7}$, against an error budget below $10^{-10}$.

\begin{remark}[$S_6$]\label{rem:tao}
At $S_6$ the horizontal spectrum of $J^\top J$ is exactly $\{9/2^{(16)},12^{(9)}\}$: an exact rational check gives $M(M-\frac92)(M-12)=0$ and the traces, and the kernel is the $11$-dimensional gauge space. So $\sigma=3/\sqrt2$ on all $25$ horizontal directions ($k=0$). Since $\Phi(0)=0$, any Hadamard $\e^{i(P+\delta)}$ with $\delta$ horizontal satisfies $\sigma\|\delta\|\le\sqrt3\|\delta\|^2$, so $\delta=0$ or $\|\delta\|\ge\sqrt{3/2}$. Thus $S_6$ is isolated. Stage C on the ball of radius $0.08$ shows that $S_6$ is not a quartet edge.
\end{remark}

\section{Step 4: the atlas}

\begin{lemma}[Boundary transfer]\label{lem:transfer}
Let $s$ and $t$ be Records, $\mathcal A$ a signed permutation of the $36$ entries (row and column permutations, transposition, conjugation) and $N$ an integer matrix. Put
\[
 c=\hor(\mathcal AP_s+2\pi N-P_t),\qquad M=\hor(\mathcal AT_s).
\]
If $\sup_{\eta\in Q}\|c+M\eta\|+n_s<R_t$ on a box $Q$, then every boundary Hadamard of $s$ with tangent coordinate in $Q$ is equivalent to an interior point of $t$.
\end{lemma}

\begin{proof}
We have $\|\hor(\mathcal A\nu)\|\le\|\nu\|\le n_s$, and representatives differ by row-plus-column patterns. Note also $\|M\|\le\|T_s\|$.
\end{proof}

The boundary therefore needs only a $4$-dimensional cover. A box is either certified off the shell $R_s-n_s\le\|T\eta\|\le R_s+n_s$, or it transfers.

\paragraph{The computation.} A breadth-first driver grows Records of radius $0.08$ until every boundary transfers. A checker recomputes every Record from its centre phases: frame, $\sigma$, normal bound, Stage C, and the boundary tree with Arb transfer constants. A second implementation, written only from the specification (Julia, outward-rounded intervals), repeats all of these checks. The atlas has \NRecords{} Records; both checkers certify every one, and $F_6$ lies in Record \Fsix.

\section{Conclusion}

By Corollary~\ref{cor:cover} and Step 4, every class in $X$ meets a Record. By Lemmas~\ref{lem:normal} and~\ref{lem:stagec}, no Hadamard in a Record is a quartet edge. By Remark~\ref{rem:tao}, neither is $S_6$. Hence $\mathbb C^6$ has no four mutually unbiased bases. \qed

\paragraph{Inputs.} The proof uses:
\begin{itemize}
\item the classification of order-6 complex Hadamard matrices [CT], through MC1--MC5, in Step 2 only;
\item the Stage C method of [OAI], with every inequality reproved;
\item certified computation in Steps 1, 3 and 4.
\end{itemize}

\paragraph{Formalisation.} The argument is formalised in Lean~4 on top of the Lean classification of order-6 Hadamard matrices. It uses only the standard axioms. Steps~1 and~3, the boundary transfer (Lemma~\ref{lem:transfer}), the open--closed atlas argument, matched uniqueness and an explicit matched completion at a maximum are proved there. The end-to-end theorem takes as hypotheses only:
\begin{itemize}
\item the finite atlas certificate;
\item Step~2, either as the connectedness of $X$ or as matched existence and density.
\end{itemize}

\begin{thebibliography}{9}
\bibitem{CT} M.~C\'ardenes Wuttig, J.~Tindall, \emph{A complete classification of complex Hadamard matrices of order six}, arXiv:2608.18053.
\bibitem{BBELTZ} I.~Bengtsson, W.~Bruzda, \AA.~Ericsson, J.-\AA.~Larsson, W.~Tadej, K.~\.Zyczkowski, \emph{Mutually unbiased bases and Hadamard matrices of order six}, J.~Math.~Phys.\ 48 (2007) 052106.
\bibitem{OAI} OpenAI, \emph{The maximum number of mutually unbiased bases in dimension six}, preprint (2026).
\end{thebibliography}
\end{document}
